% Chapter 1. Each chapter is its own file; add new ones in main.tex.
\chapter{Introduction}
\label{ch:introduction}

\section{Storage that nobody can weigh}
\label{sec:weigh}

Between 2019 and 2025 the installed capacity of grid-connected battery
storage in north-western Europe grew more than tenfold
\parencite{iea2025storage}. Most of these systems earn their keep by
providing frequency containment reserve: they absorb or inject power
within seconds whenever the grid frequency drifts from its nominal
50\,Hz, and they are paid for being available rather than for the energy
they move. A typical site is a row of shipping containers, each holding
several thousand cylindrical or prismatic cells, a battery management
system and an inverter.

Every cell in such a site loses capacity as it ages. The loss is slow,
a few per cent a year, but it matters twice over. It erodes the energy
that the operator can bid into the market, and it determines when the
site must be refurbished, a decision worth several million euros. Yet
the quantity that governs both, the state of health of the battery, is
not measured during operation. It can only be inferred.

\begin{definition}{State of health}{soh}
  The \emph{state of health} of a battery at time \(t\) is the ratio of
  its present usable capacity to its capacity at commissioning,
  \begin{equation}
    \mathrm{SoH}(t) = \frac{Q(t)}{Q_0},
    \label{eq:soh}
  \end{equation}
  where both capacities are measured by a full discharge at a reference
  current and temperature. A battery is conventionally said to reach its
  end of life when \(\mathrm{SoH}\) falls below 0.8.
\end{definition}

The definition is simple, but measuring \(Q(t)\) requires taking the
site out of service, cycling it fully at low current and putting it back,
a reference performance test that costs a day of revenue and is therefore
performed at most once or twice a year. Between tests, operators rely on
the estimate reported by the battery management system, which is
typically a coulomb counter corrected by a lookup table supplied by the
cell manufacturer. Several audits have found these estimates to drift by
five per cent or more within a year \parencite{halvorsen2023bms}.

\section{Why laboratory models fall short}
\label{sec:lab}

Battery ageing has been studied for three decades, and the mechanisms are
well understood at the level of a single cell
\parencite{birkl2017degradation}. Lithium is consumed by the slow growth
of the solid electrolyte interphase on the anode; active material
cracks and loses contact under the mechanical stress of repeated
lithiation; and at low temperature or high current, metallic lithium can
plate on the anode surface. Each mechanism has a characteristic
dependence on temperature, state of charge and current, which
semi-empirical models capture with a handful of fitted parameters
\parencite{schmalstieg2014holistic}.

The difficulty is one of transfer. Laboratory campaigns cycle small
cells at fixed conditions chosen to separate the mechanisms, while a
frequency-response site holds most of its cells near half charge and
dithers around that point in shallow, irregular micro-cycles. Cells in
the middle of a rack run several degrees warmer than those at its edge.
None of these conditions is represented in the data on which the
parameters were fitted, and the error that results is not random but
systematic: it grows with time.

\begin{figure}[tb]
  \marginfigure{%
  \begin{tikzpicture}[
      x=1mm, y=1mm, font=\sffamily\scriptsize,
      block/.style={draw=accent, line width=0.6pt, fill=accentlight,
        minimum width=27mm, minimum height=10mm, align=center},
      plain/.style={draw=rulegrey, line width=0.4pt,
        minimum width=27mm, minimum height=10mm, align=center},
      flow/.style={-{Stealth[length=1.8mm]}, line width=0.5pt, rulegrey}]
    \node[plain] (tele) at (0,24) {Site telemetry\\\textcolor{inkmuted}{1 Hz, 3 years}};
    \node[plain] (rpt)  at (0,6)  {Reference tests\\\textcolor{inkmuted}{twice a year}};
    \node[block] (phys) at (46,24) {Ageing law\\\textcolor{inkmuted}{Chapter 2}};
    \node[block] (gp)   at (46,6)  {GP correction\\\textcolor{inkmuted}{Chapter 3}};
    \node[plain] (soh)  at (92,15) {State of health\\\textcolor{inkmuted}{with uncertainty}};
    \draw[flow] (tele) -- (phys);
    \draw[flow] (tele.east) ++(0,-2) -- (gp.north west);
    \draw[flow] (rpt) -- (gp);
    \draw[flow] (phys) -- (gp);
    \draw[flow] (phys.east) -- (soh.north west);
    \draw[flow] (gp.east) -- (soh.south west);
  \end{tikzpicture}}{%
  \caption[Structure of the hybrid estimator]{Structure of the hybrid
    estimator. Operational telemetry drives a physical ageing law; a
    Gaussian process learns its residual from sparse reference tests.}
  \label{fig:overview}}
\end{figure}

\section{Approach and contributions}
\label{sec:contributions}

This thesis takes the opposite route. Rather than transferring a
laboratory model to the field, it learns from the field directly and uses
physics only where data are silent. \Cref{fig:overview} shows the
resulting structure, which the remaining chapters build up in turn. The
work makes three contributions:

\begin{itemize}
  \item an open, harmonised data set of operational telemetry and
    reference tests from eleven grid-scale sites, described in
    \cref{ch:data};
  \item evidence that calendar ageing at elevated state of charge, not
    cycle depth, dominates capacity fade under frequency-response duty
    (\cref{sec:drivers});
  \item a hybrid estimator with calibrated uncertainty, and a field trial
    showing that it can safely enlarge the usable state-of-charge window
    (\cref{ch:estimation}).
\end{itemize}

The scope is deliberately narrow. All eleven sites use one of two cell
chemistries in air-cooled containers and serve the same frequency-response
market, so the findings say little about liquid-cooled racks or
arbitrage-heavy duty. Within that scope, however, the sites differ in
climate, controller settings and operator practice enough to test whether
a model trained on some of them carries over to the others, which is the
question an operator commissioning a new site actually faces.

Parts of this work have been published before. The data set and the
analysis of ageing drivers appeared in two journal articles, and an early
version of the estimator was presented at a conference on grid storage.
Each chapter notes where it extends or corrects that earlier material.

The thesis is organised in two parts. Part~I sets out the problem and the
data. Part~II presents the estimator, its evaluation and the field trial,
and closes with the limits of what operational data can reveal.
\Cref{app:protocol} documents the reference test protocol in enough
detail to reproduce it.
